% Generated draft: requires mathematical and source review.
\documentclass[11pt]{article}
\usepackage[margin=25mm]{geometry}
\usepackage{fontspec}
\setmainfont{Latin Modern Roman}
\usepackage{amsmath,amssymb,mathtools,unicode-math}
\setmathfont{Latin Modern Math}
\usepackage{enumitem,needspace,xurl,hyperref}
\setlength{\parindent}{0pt}
\setlength{\parskip}{5pt}
\setlength{\emergencystretch}{3em}
\newcommand{\N}{\mathbb{N}}
\newcommand{\Z}{\mathbb{Z}}
\newcommand{\Q}{\mathbb{Q}}
\newcommand{\R}{\mathbb{R}}
\newcommand{\C}{\mathbb{C}}
\newcommand{\F}{\mathbb{F}}
\newcommand{\A}{\mathbb{A}}
\newcommand{\PP}{\mathbb{P}}
\newcommand{\E}{\mathbb{E}}
\newcommand{\eps}{\varepsilon}
\newcommand{\dd}{\,\mathrm{d}}
\begin{document}
\section*{UnsolvedMath ID: 10400121}
\textbf{Source problem code:} catalogue entry
\subsection{Definitions and mathematical statement}
For the source's compact gauge-group/level conventions, let $M$ be a closed oriented three-manifold and $Z_k^G(M)$ its quantum invariant. Put $r=k+h^\vee$, with the source's dual Coxeter shift. The perturbative expansion conjecture proposes the following formal large-level structure:
\[
Z_k^G(M)\sim e^{-\pi i(\dim G)(1+b_1(M))/4}
\int_{[A]\in\mathcal M_G(M)}
e^{2\pi i r\operatorname{CS}(A)}
r^{(h_A^1-h_A^0)/2}
 e^{-2\pi i(I_A/4+(h_A^0+h_A^1)/8)}
\tau_M(A)^{1/2}
\exp\left(\sum_{\ell\geq1}a_\ell(k)
\sum_{e(\Gamma)=-\ell}\frac{Z_\Gamma(M,A)}{|\operatorname{Aut}(\Gamma)|}\right).
\]
Here $\mathcal M_G(M)$ is the intended flat-connection moduli space, $\operatorname{CS}$ the normalized Chern--Simons phase, and $h_A^j$ naturally denotes twisted-cohomology dimensions $\dim H^j(M;\operatorname{ad}A)$. The spectral-phase quantity $I_A$, torsion factor $\tau_M(A)$ and graph integrals $Z_\Gamma$ require the source's definitions.

The symbol $a_\ell(k)$ in this display only stands for the extracted coefficient involving the printed $c\ell$, $k^{-\ell}$, $(2\ell)!$ and $(3\ell)!$. Their superscript/subscript and fraction layout is not recoverable unambiguously from the flattened record; no factorial ratio is guessed. Likewise $e(\Gamma)=-\ell$ is retained as source graph notation, not interpreted as a negative number of edges. It may encode an Euler-characteristic convention that must be checked.

The task is to give the right side a mathematical interpretation and establish the asymptotic relation in the intended generality. The source explicitly says it is mathematically defined in certain preceding cases but needs further interpretation in general. The measure/density on possibly singular or positive-dimensional moduli, stabilizer treatment, phase/square-root conventions, graph class and regularization, group/level range, and meaning of the asymptotic remainder are not supplied. In nonacyclic situations torsion need not be a scalar without additional data. A formal perturbative series alone does not prove an asymptotic expansion of the numerical quantum invariant. No universal ordinary integral or already-proved general formula is asserted.
\subsection{Short English statement}
Make rigorous the proposed flat-connection/graph expansion of large-level quantum invariants, with the exact phases, cohomological powers and dual-Coxeter shift. The source explicitly leaves general moduli-space interpretation and some extracted coefficient notation unresolved.
\subsection{Sources}
\begin{itemize}
\item Local UnsolvedMath registry, permanent ID 10400121, exact source code catalogue entry, Conjecture 7.6, ``The perturbative expansion conjecture''; no proposer/year supplied. Registry status partially solved, classification PARTIAL-PROGRESS.
\item Ohtsuki, \emph{Problems on invariants of knots and 3-manifolds} (2002), Conjecture 7.6, PDF page 105, \url{https://msp.org/gtm/2002/04/gtm-2002-04-024s.pdf}. Supplied pdf-font-structure extraction accessed July 29, 2026; provenance distinguishes 231 bold numbered items. The preceding special cases and definitions are absent; no PDF was accessed here.
\item The supplied summary says perturbative Chern--Simons theory gives formal understanding but supplies no precise rigorous all-manifold expansion theorem. Formal progress is not identified with the conjectured numerical asymptotics.
\item Provenance flags current primary-literature status and redistribution terms NEEDS\_REVIEW.
\end{itemize}
\end{document}
